% 第四部分：固定实验记录的同题比较。
\FloatBarrier
\Needspace{6\baselineskip}
\begingroup
\setlength{\parskip}{3bp}
\setlength{\emergencystretch}{1em}
\clubpenalty=10000
\widowpenalty=10000
\setcounter{figure}{0}
\setcounter{table}{0}
\renewcommand{\thefigure}{4-\arabic{figure}}
\renewcommand{\thetable}{4-\arabic{table}}
\renewcommand{\theHfigure}{section4.\arabic{figure}}
\renewcommand{\theHtable}{section4.\arabic{table}}
\renewcommand{\figurename}{图}
\renewcommand{\tablename}{表}
\captionsetup{font=reportcaption,labelsep=quad,skip=5bp}
\newcommand{\ChFourFigure}[4][0.98]{%
  \begin{figure}[!htbp]\centering
  \includegraphics[width=#1\linewidth]{section4Figures/#2}
  \caption{#3}\label{#4}\end{figure}}
\newcommand{\ChFourTable}[4]{%
  \begin{table}[!htbp]\centering\caption{#2}\label{#3}
  \begingroup\fontsize{9bp}{12bp}\selectfont
  \setlength{\tabcolsep}{3bp}\renewcommand{\arraystretch}{1.15}
  \input{section4Tables/#1.tex}\par\vspace{3bp}
  \begin{minipage}{0.98\linewidth}\fontsize{8.5bp}{11.5bp}\selectfont #4\end{minipage}
  \endgroup\end{table}}

\TemplateSection{4.\quad N-Puzzle 问题求解探讨}
\noindent{\bfseries\fontsize{12bp}{15.6bp}\selectfont 4.\quad N-Puzzle 问题求解探讨}\par\vspace{6bp}

本章固定初态、目标和单位移动代价，比较 A* 与 IDA* 的实际表现，
并考察启发估值、转置表和快速后继生成的收益与代价。
结点的评价函数为 $f(n)=g(n)+h(n)$：$g$ 是已经移动的步数，$h$ 是剩余步数下界。
更大的下界可能减少扩展，但运行时间还取决于估值、对象构造、查重和数据库准备的成本。
因此，实验同时报告成功次数、时间和扩展事件数，不把其中任一项单独当作优劣标准。

\textbf{实验对象与配置。}
教师三阶题使用第 2.1 节的初态与非标准目标，解长为 16；
普通四阶题取 \texttt{problems3.txt} 第 4 题，解长为 52；
困难四阶题取同文件第 10 题，即教师指定的倒序棋盘，成功解长为 70。
所有配置使用 JDK 19.0.2，源码以 \texttt{--release 19} 编译；系统为 Windows，处理器为 AMD Ryzen 9 8940HX，最大堆为 \texttt{-Xmx1024m}。
每次启动独立 JVM，不在同一 JVM 内预热。
所有运行均设置 60 秒上限，从启动求解进程前开始计时。

三阶和普通四阶分别对两种算法、四种启发函数各重复 3 次，共 48 次。
困难题比较 8 种模式与开关配置，各重复 3 次，共 24 次，另做 1 次独立冷缓存运行。
加上 7 次用于检查配置和建立私有缓存的预实验，计划共 80 次，实际为 3 次 OOM、74 次 SOLVED、3 次 TIMEOUT。
预实验与独立对照另行统计。三阶各配置均成功，无超时；普通四阶错位数中，A* 三次均因堆内存不足（OOM）结束，IDA* 三次均达到 60 秒上限（TIMEOUT）。

\textbf{计时和统计口径。}
搜索时间从调用搜索器到返回路径，包括启发查询和运行器的同步快照记录，
不含提前构建或加载 PDB。求解总耗时从启动求解进程前计至进程退出或超时终止回收，
包含 JVM 启动、输入、准备、搜索、路径回放和导出。两种计时分别反映搜索本身与一次求解的整体成本。
表内时间分别取成功运行的中位数，失败不填充为上限值，也不参加成功中位数；
同时列出成功次数和失败类型。每种配置重复三次，结果用于描述本组固定题目的性能。

PDB 在三阶使用 4-4 分组，普通四阶使用 4-4-4-3，困难题使用 555T。
四种启发函数对照中的 PDB 每次重新建表；困难题模式对照复用预实验建立的同一份缓存。独立冷运行单列建表与总耗时，热运行单列缓存加载与总耗时。
图~\ref{fig:ch4-protocol} 示意这些计时范围。
\ChFourFigure{unified/protocol.pdf}{求解进程的执行步骤与计时范围}{fig:ch4-protocol}
\FloatBarrier

\TemplateSubsection[templatetwoextra]{4.1 Open 表及 Closed 表数据结构的改进及其对算法性能的影响}
\noindent{\bfseries 4.1\quad Open 表及 Closed 表数据结构的改进及其对算法性能的影响}\par\vspace{2bp}

A* 的 Open 需要同时支持选出最小 $f$ 和定位重复状态。
\texttt{PqFrontier} 用优先队列保存候选结点，再用状态索引维护每个棋盘的有效记录。
堆的入队、出队操作为 $O(\log N)$，哈希索引查询平均为 $O(1)$。
替换旧结点时，\texttt{queue.remove(oldNode)} 需要线性定位，因此替换操作为 $O(N)$。
有序列表即使用二分查找确定位置，插入与删除仍可能搬移元素。表~\ref{tab:ch4-operations} 列出结构代价。
\ChFourTable{operations}{Open 与 Closed 的结构代价}{tab:ch4-operations}{%
$N$ 为 Open 中的结点数。这些是操作复杂度，不是实测耗时或内存峰值。}

Closed 使用 \texttt{HashSet<State>}，按完整状态判等。
本实验固定 A* 为优先队列与状态索引的组合，Closed 也保持相同实现，
从而使不同启发函数使用相同的候选管理规则。

这套结构将频繁的最小值出队与状态查询交给堆和哈希索引，
额外代价是索引存储，以及更新较优重复结点时的线性删除。
本章报告整套实现的运行时间，未设置列表版的配对计时。
\ChFourFigure{unified/open-operations.pdf}{候选后继的状态定位、插入与替换}{fig:ch4-open}
\FloatBarrier

\TemplateSubsection{4.2 A* 算法、IDA* 算法的性能比较}
\noindent{\bfseries 4.2\quad A* 算法、IDA* 算法的性能比较}\par\vspace{2bp}

两种算法使用同一个 $f=g+h$，但保留的信息不同。A* 同时保存 Open、状态索引、Closed 和父链；
BASIC IDA* 在阈值内深度优先搜索，主要保留当前路径，升高阈值后从根重新搜索。这使第二阶段能够把搜索空间从三阶扩展到四阶：以跨轮重复计算换取较小的路径存储，但启发信息过弱时仍可能耗尽时间。四阶错位数的 A* 内存不足与 IDA* 超时，分别体现了两种资源限制。
本节 IDA* 采用 BASIC、\texttt{fast=true}，不启用转置表。

表~\ref{tab:ch4-teacher3} 给出教师三阶题的全部配置，3 次重复均成功，解长均为 16。
同用 Manhattan 时，A* 搜索中位数为 6.002\,ms，
IDA* 为 18.828\,ms，扩展事件数分别为 311 和 423。
A* Manhattan 的求解总耗时中位数为 145.389\,ms，
最大值为 170.231\,ms，
该配置 3 次都在 1 秒内完成，满足第一阶段要求。
\ChFourTable{current-teacher3}{教师三阶题的算法与启发函数比较}{tab:ch4-teacher3}{%
每配置 3 次，均为 16 步。时间为各指标分别计算的成功中位数；PDB 为冷构建 4-4。
LC 指 Manhattan 加线性冲突修正。建表时间不包含在搜索时间内。}

普通四阶题的结果见表~\ref{tab:ch4-ordinary4}。
Manhattan 下，A* 搜索中位数为 685.852\,ms，
IDA* 为 196.860\,ms；扩展分别为 339\,410 与 997\,736。
IDA* 的扩展包含跨轮重复，而 A* 需要维护全局候选集合。
在该题上，IDA* 扩展更多，但搜索耗时更短，反映了两种实现的单次扩展成本差异。
错位数的失败状态单独保留，其余成功配置的路径均为 52 步。
\ChFourTable{current-ordinary4}{普通四阶题的算法与启发函数比较}{tab:ch4-ordinary4}{%
每配置 3 次，普通题上限 60 秒。TO 为 TIMEOUT，OOM 为内存不足；失败不填完成指标，成功次数与失败类型同时保留；
其余配置 3/3 成功、解长 52，PDB 冷构建 4-4-4-3。}

生成事件和扩展事件必须结合代码理解。IDA* 的计数跨轮累加，
fast 还会在创建子状态前过滤立即反向动作；A* 则保留前沿和 Closed。
表中的生成与扩展均为累计事件数，包含同一状态的重复访问。
存储分析按算法保存的数据结构展开；\texttt{-Xmx1024m} 为堆上限，本实验未测量 JVM 峰值内存。
\ChFourFigure{unified/search-storage.pdf}{A* 与 BASIC IDA* 保留的信息；框面积不表示实际内存}{fig:ch4-storage}
\FloatBarrier

\TemplateSubsection{4.3 不同启发函数下 A* 与 IDA* 的性能比较}
\noindent{\bfseries 4.3\quad 不同启发函数下 A* 与 IDA* 的性能比较}\par\vspace{2bp}

课程要求比较错位数、Manhattan 与加性 PDB。本章在相同状态对上同时列出 LC，
用来检查更强估值能否抵消额外计算成本。错位数只区分是否归位，Manhattan 保留位移，
LC 进一步约束同一目标行列中将牌的顺序，而 PDB 保存各组的抽象最短距离。

PDB 可加的关键是成本划分：一次动作只由被移动将牌所在组收费 1，其余组收费 0。
抽象距离通过 0-1 BFS 建立。不同完整分组之间有重叠时不能直接相加；
各自为有效下界时可以取最大值。图~\ref{fig:ch4-pdb} 给出准备和查表关系。
\ChFourFigure{unified/pdb-pipeline.pdf}{加性 PDB 的准备、校验与估值流程}{fig:ch4-pdb}
\FloatBarrier

\TemplateSubsubsection{4.3.1 运行时间对比及分析}
\noindent{\bfseries 4.3.1\quad 运行时间对比及分析}\par\vspace{2bp}

三阶 A* 的 Manhattan 与 LC 搜索中位数分别为 6.002、
5.061\,ms；IDA* 分别为 18.828、
19.593\,ms。
普通四阶 A* 的对应搜索中位数为 685.852、508.963\,ms，
IDA* 为 196.860、364.377\,ms。
LC 减少扩展，但估值本身也要计算行列顺序；时间是否下降应分别按算法和题目判断。
普通四阶题中，LC 使 A* 的搜索更快，却使 IDA* 的搜索更慢。
LC 减少扩展的同时也增加了行列顺序的计算，IDA* 的重复估值进一步累积了这部分成本。
\ChFourFigure{current/heuristic-search.pdf}{相同题目下四种启发函数的搜索时间中位数}{fig:ch4-times}

计入准备成本后，三阶 PDB 的 A* 与 IDA* 求解总耗时中位数分别为 219.227、234.436\,ms；
普通四阶的对应值分别为 743.599、585.805\,ms。
这些值包含每次冷建表的成本，必须与表内同算法、同题的 Manhattan 求解总耗时比较。
若多次复用同一目标和分组，应把首次建表与每次加载分别计入。
普通四阶中，A* 使用 PDB 时，搜索与总耗时均低于 Manhattan；
IDA* 使用 PDB 时，搜索更短，总耗时却因建表而更长。
因此，普通题的一次性求解需要权衡建表与搜索成本，困难题则更依赖 PDB 提供的强下界。
\ChFourFigure{current/heuristic-parent.pdf}{包含准备、搜索、回放和导出的求解总耗时中位数}{fig:ch4-parent}
\FloatBarrier

\TemplateSubsubsection{4.3.2 扩展结点比较及分析}
\noindent{\bfseries 4.3.2\quad 扩展结点比较及分析}\par\vspace{2bp}

普通四阶 A* 的扩展事件由 Manhattan 的 339\,410 降到 LC 的 184\,819，
IDA* 由 997\,736 降到 595\,024。两种算法都少扩展，但时间变化方向不同，说明估值强度不是唯一成本。
PDB 在该题进一步把 A* 扩展减至 118\,544、IDA* 减至 275\,777；
它的建表和查表开销仍需通过表~\ref{tab:ch4-ordinary4} 的时间列一起判断。
\ChFourFigure{current/lc-nodes.pdf}{普通四阶题 Manhattan 与 LC 的扩展事件数}{fig:ch4-expanded}

困难题 BASIC 成功运行的阈值依次为 $56,58,60,62,64,66,68,70$，在第 8 轮找到 70 步解。
阈值由最小超界值推进，本题形成间隔为 2 的序列。
成功运行累计扩展 48\,759\,798 次，包含跨轮重复访问。
\FloatBarrier

\TemplateSubsection{4.4 转置表模式、容量与缓存}
\noindent{\bfseries 4.4\quad 转置表模式、容量与缓存}\par\vspace{2bp}

困难题初态为 $[15,14,13,12;\ 11,10,9,8;\ 7,6,5,4;\ 3,1,2,0]$，
目标为 $1$ 至 $15$ 顺序排列、右下角为空位。555T 的分组为
\[
P_1=\{1,5,9,13,2\},\quad P_2=\{6,10,14,3,7\},\quad P_3=\{11,15,4,8,12\}.
\]
困难题使用 555T，普通四阶使用 4-4-4-3，两组结果分别统计。
模式对照复用同一份缓存，每次搜索启动独立 JVM。
TT 与 ZOBRIST 的容量分别设为 10 万、50 万和 100 万，达到上限后不接纳新状态，
但已有状态发现较小 $g$ 时仍可更新。
\ChFourTable{current-modes}{困难题热缓存下的模式与容量比较}{tab:ch4-modes}{%
每配置 3 次，本组 24 次均成功，解长均为 70。
时间是成功样本中位数，需与成功次数一起阅读。容量 0 的 BASIC 不使用转置表。}

在 3 次均成功的配置中，求解总耗时中位数最低者为 BASIC、容量 0、fast 开启，
为 13.791 秒。表中同时保留其他配置的成功次数与耗时。
转置表可能减少扩展，却增加哈希、完整判等、分配和记录成本；
TT 与 ZOBRIST 的三档容量均完成求解，但总耗时均高于 BASIC、fast 的组合。
增加容量并未使总耗时持续下降。因此，本实验选用 BASIC、fast 作为该困难题的配置，
并保留转置表用于其他重复状态分布下的求解。

独立冷启动采用 BASIC、fast、555T，从头构建模式数据库。
该次建表 2985\,ms，搜索 13\,816.796\,ms，
求解总耗时 16\,981.279\,ms，状态为 SOLVED，得到 70 步合法路径。
这次独立冷启动在 60 秒内完成，满足困难题的 1 分钟要求。
表~\ref{tab:ch4-hard} 分别列出独立冷启动与三次热缓存运行的结果。
\ChFourTable{current-cache}{困难题 BASIC、fast、555T 的缓存条件与结果}{tab:ch4-hard}{%
冷缓存为独立 1 次、70 步；热缓存为模式对照中的 3 次，各时间列分别取中位数。
冷、热两种条件分别统计。}
\FloatBarrier

\TemplateSubsection{4.5 fast 分支与记录开销}
\noindent{\bfseries 4.5\quad fast 分支与记录开销}\par\vspace{2bp}

fast 在状态构造前排除立即反向动作，随后提前检查路径环，并逐个创建子结点。
Manhattan 和错位数还可通过增量接口更新单牌贡献；PDB 没有该接口。
这些机制由同一个开关控制，本组对照测得其组合效果。

困难题关闭 fast 时为 3/3 成功，
开启时为 3/3。
第 3 组完整配对的搜索时间分别为 19.804 与 11.521 秒。
表~\ref{tab:ch4-counts} 给出这一组的工作量：生成事件减少
48\,759\,757，恰好等于提前过滤的反向动作数。
状态、结点构造及 PDB 查询还受到提前路径检查、逐个生成和未访问兄弟结点的影响。
两边扩展次数相同，这组差异主要反映每次扩展所做工作的变化；
表中列出第 3 组的具体计数，时间比较采用三次运行结果。
\ChFourTable{current-counts}{困难题第3组完整配对的构造与查询计数}{tab:ch4-counts}{%
同为 BASIC、555T、热缓存、70 步。计数是实际调用或事件累计，不是对象存活数或内存峰值。}

独立普通题控制使用相同的 52 步实例、BASIC、555T 热缓存，交叉切换 fast 与快照开关，
每种组合在独立 JVM 中重复运行三次，按固定种子打乱运行顺序。关闭快照时，fast 关闭与开启的搜索中位数分别为
82.132 和 62.117\,ms；
开启快照后，对应为 94.814 和 75.798\,ms。
同步记录属于搜索区间内的额外工作，比较时必须保持相同记录条件。
该对照通过反射调用计数导出接口，包含少量诊断开销，结果用于比较普通题中的快照开关影响。
\ChFourTable{current-controls}{普通四阶题的fast 与快照交叉控制}{tab:ch4-controls}{%
每种组合 3/3 成功、解长 52；使用 555T 热缓存，不与表~\ref{tab:ch4-ordinary4} 的 4-4-4-3 冷构建混合。
状态构造数与 PDB 查询数按实际接口调用统计。}
\FloatBarrier

\TemplateSubsection{4.6 正确性核验与适用范围}
\noindent{\bfseries 4.6\quad 正确性核验与适用范围}\par\vspace{2bp}

独立棋盘编码与相邻交换 BFS 对三个目标各枚举 181\,440 个可达三阶状态，
检查 LC 的目标零值、Manhattan 下界、精确距离上界和每条边的一致性。
四阶另用浅层 BFS、48 组 A* Manhattan 精确解和随机邻边核对。
转置表测试检查强制碰撞、容量边界、较小 $g$ 更新、逐轮清空和阈值推进，
并在 300 个带环随机图上对照独立最短距离。专项共 7\,022\,745 项检查，9 套 Java 测试记录均通过。
四阶与随机图的核对覆盖所选实例。

所有成功求解的记录均逐步回放，检查初态、动作连续性、目标状态和解长。
超时与内存不足单列在状态统计中。各配置的运行环境和原始数据见仓库的数据索引。

实验结果表明，启发估值、候选管理和准备成本共同影响运行时间。
LC 减少了扩展次数，但在 A* 与 IDA* 中产生了不同的时间效果；
PDB 将部分工作移到建表阶段，需结合冷、热缓存分别评价。
第一阶段 A* 的 Manhattan 配置满足 1 秒要求；困难题独立冷启动得到 70 步解，包含建表的总时间为 16.981 秒，满足 1 分钟要求。
解路径文件的可视化展示见第五部分。
\FloatBarrier
\endgroup
\par\vspace{8bp}
\Needspace{6\baselineskip}
